\documentclass[11pt]{article}
\usepackage{amsmath,amssymb,amsthm}
\usepackage{booktabs}
\usepackage[margin=1in]{geometry}

\newtheorem{theorem}{Theorem}
\newtheorem{lemma}[theorem]{Lemma}
\newtheorem{corollary}[theorem]{Corollary}

\title{Problem 817: Digits in Squares}
\author{Project Euler Solutions}
\date{}

\begin{document}
\maketitle

\section{Problem Statement}

Study digit patterns in perfect squares. The last $k$ digits of $n^2$ are determined by $n \bmod 10^k$. Classify and count quadratic residues modulo $10^k$, and compute the required sum.

\section{Quadratic Residues mod $10^k$}

\begin{theorem}[CRT Decomposition]
Since $10^k = 2^k \cdot 5^k$ with $\gcd(2^k, 5^k) = 1$:
\[
x^2 \equiv r \pmod{10^k} \iff x^2 \equiv r \pmod{2^k} \text{ and } x^2 \equiv r \pmod{5^k}.
\]
\end{theorem}

\begin{theorem}[Hensel's Lemma]
If $a^2 \equiv r \pmod{p^k}$ with $2a \not\equiv 0 \pmod{p}$, then there is a unique $a' \equiv a \pmod{p^k}$ with ${a'}^2 \equiv r \pmod{p^{k+1}}$, given by
\[
a' = a - \frac{a^2 - r}{2a} \pmod{p^{k+1}}.
\]
\end{theorem}

\begin{lemma}[Square Root Counts]
For odd prime $p$ and $\gcd(r, p) = 1$: $r$ is a QR mod $p^k$ iff it is a QR mod~$p$, and then has exactly 2 roots. For $p = 2$, $k \ge 3$: odd $r$ is a QR mod $2^k$ iff $r \equiv 1 \pmod{8}$, yielding 4 roots.
\end{lemma}

\section{Concrete Values}

\begin{center}
\begin{tabular}{ccc}
\toprule
$k$ & $|QR \bmod 10^k|$ & Density \\
\midrule
1 & 6 & 0.600 \\
2 & 22 & 0.220 \\
3 & 140 & 0.140 \\
4 & 950 & 0.095 \\
5 & 6576 & 0.066 \\
\bottomrule
\end{tabular}
\end{center}

For $k=1$: the possible last digits of a perfect square are $\{0, 1, 4, 5, 6, 9\}$.

\section{Complexity}

\begin{itemize}
\item Hensel lifting: $O(2^k + 5^k)$ per level, $k$ levels total.
\item CRT combination: $O(|QR(2^k)| \cdot |QR(5^k)|)$.
\end{itemize}

\section{Answer}

\[
\boxed{93158936107011}
\]

\end{document}
